% Ad Familiares 14.19 --- headnote
% ad-familiares-14-19, 28 November 48 BC, written at Brundisium

\textit{Cicero to Terentia, written from Brundisium on the
fourth day before the Kalends of December 48 BC (the
manuscript dateline: \textit{Scr. Brundisi iv K. Dec.
a. 706 (48)}). Three weeks after the homecoming note of
\textit{Fam.} 14.12, Cicero is still stuck at Brundisium
under the same shadow; Tullia, now married to Dolabella,
is unwell at Rome, and Terentia has been urging him to
come up the road toward her.}

\textit{The letter is six lines long and registers three
things in turn: dread for Tullia's health (he will not
elaborate, knowing it weighs on Terentia as heavily as on
him); the admission that he ought to be moving closer to
them and would have already, but for impediments not yet
cleared; and a courier instruction --- he is waiting on a
letter from Pomponius (Atticus), and Terentia is to make
sure it reaches him at speed. The closing ``take pains to
keep well'' is the same formulation he has been writing to
her all autumn.}
